\clearpage
\SourceLeaf{402}\pdfbookmark[1]{II. Tafeln für die Störungen der Ceres}{ceres-tables-letters}
\begin{center}
{\large BRIEFWECHSEL.\par}
\bigskip
[II. TAFELN FÜR DIE STÖRUNGEN DER CERES.]\par
\bigskip
\textsc{Gauss} an \textsc{Olbers}, Braunschweig, 3. Dezember 1802.
\end{center}
\noindent\hspace*{\parindent}.\quad.\quad.\quad.\quad.\quad Ich berechne mir jetzt eine Tafel für die im \rlap{\kern.02em\raisebox{1.02em}{\rule{.2em}{.25pt}}}Nov. der M. C. a\rlap{\kern.35em\raisebox{1.02em}{\rule{.2em}{.25pt}}}ngezeigten Störungen der {\AstroSymbols ⚳}. Dabei weiche ich in der Einrichtung zum Teil vom üblichen Verfahren ab. Die Tafeln für die nur von \(\text{\AstroSymbols ⚳}-\text{\AstroSymbols ♃}\) abhängigen Störungen der Länge und des Radiusvektors und die für die Breite sind schon fertig. Ich schreibe Ihnen darüber demnächst ausführlicher.\quad.\quad.\quad.\quad.\quad.

\begin{center}
\textsc{Gauss} an \textsc{Olbers}, Braunschweig, 21. Dezember 1802.
\end{center}
\noindent\hspace*{\parindent}.\quad.\quad.\quad.\quad.\quad Meine jetzt fertiggestellten Tafeln\refstepcounter{footnote}\hyperlink{ceresII-note}{[*)]}\begingroup
\makeatletter\long\def\@makefntext#1{\noindent[*) #1}\makeatother
\footnotetext{\hypertarget{ceresII-note}{}Band VI, Seite 235—243.]}
\endgroup{} für die Störungen der {\AstroSymbols ⚳} durch {\AstroSymbols ♃} werde ich noch in dieser Woche an \textsc{Zach} schicken. Sind sie für die M. C. noch zu umfangreich, so werde ich Ihnen eine Abschrift besorgen. Ihnen kann ich die Einrichtung mit ein paar Worten beschreiben. Bei der Länge bilden zunächst alle Gleichungen, deren Argument ein Vielfaches von \(\text{\AstroSymbols ⚳}-\text{\AstroSymbols ♃}\) ist, eine Tafel. Die Summe aller übrigen Gleichungen, die von der einfachen Exzentrizität abhängen, lässt sich durch \(A\sin(B-\text{\AstroSymbols ⚳})\) ausdrücken, sodass \(A\) und \(B\) nur Funktionen von \(\text{\AstroSymbols ⚳}-\text{\AstroSymbols ♃}\) sind. Diese Größen \(A\) und \(B\) gibt eine zweite Tafel für alle Werte von \(\text{\AstroSymbols ⚳}-\text{\AstroSymbols ♃}\) von Grad zu Grad an. Ebenso kommen 2 Tafeln für den Radiusvektor und eine für die Breite hinzu, zusammen 5 Tafeln, da man sonst für alle von mir verwendeten Gleichungen 40 Tafeln oder doch an die 30 brauchen \SourceLeaf{403}würde, wenn man diejenigen wegließe, die nicht auf \(2''\) kommen. Mir wenigstens ist es \textls[180]{weit} leichter, diese einfachen trigonometrischen Rechnungen auszuführen, als so viele Argumente zu bilden, in so vielen Tafeln nachzuschlagen und so viele Additionen auszuführen. Man könnte den Gebrauch meiner Tafel leicht noch weiter erleichtern, wenn man zum Beispiel statt \(A\) gleich \(\log A\) aufnähme. Man könnte auch nach Belieben eine konstante Größe \(=2M\) wählen, die mindestens so groß wäre wie jeder Wert von \(A\), und \(\frac{A}{M}=\cos C\) setzen. Dadurch würde die Summe der von der Exzentrizität abhängigen Gleichungen
\[
=M\sin(B-C-\text{\AstroSymbols ⚳})+M\sin(B+C-\text{\AstroSymbols ⚳}).
\]
Man nähme dann in die Tafel für jeden Wert von \(\text{\AstroSymbols ⚳}-\text{\AstroSymbols ♃}\) die Werte von \(B-C\) und \(B+C\) auf und fügte noch eine Tafel hinzu, die \(M\sin\varphi\) für alle Werte von \(\varphi\) enthielte. — In meiner Tafel habe ich übrigens anstelle der Winkel \(B\) für
\[
\begin{array}{@{}l@{\quad}l@{}}
\text{die Länge}\quad\text{.\ .\ .\ .}&B-2(\text{\AstroSymbols ⚳}-\text{\AstroSymbols ♃})\\[3pt]
\left.\begin{array}{@{}l@{}}\text{den R. V.}\\\text{Breite}\end{array}\right\}\ \text{.\ .\ .}&B-3(\text{\AstroSymbols ⚳}-\text{\AstroSymbols ♃})
\end{array}
\]
die angegebenen Werte angesetzt oder der Summe der Gleichungen die Form
\[
\begin{array}{@{}l@{\quad}l@{}}
\text{für die Länge}\ \text{.\ .\ .\ .\ .\ .\ .\ .\ .\ .\ .\ .\ .\ .\ .}&A\sin(B'+\text{\AstroSymbols ⚳}-2\text{\AstroSymbols ♃})\\[4pt]
\text{für den R. V. und die Br.}\ \text{.\ .\ .\ .}&A\sin(B'+2\text{\AstroSymbols ⚳}-3\text{\AstroSymbols ♃})
\end{array}
\]
gegeben.

Dieses \(B'\) bleibt so stets innerhalb bestimmter Grenzen, wenn \(\text{\AstroSymbols ⚳}-\text{\AstroSymbols ♃}\) alle Werte durchläuft, während \(B\) in einer Periode von \(\text{\AstroSymbols ⚳}-\text{\AstroSymbols ♃}\) einen oder zwei Umläufe machen würde.\quad.\quad.\quad.\quad.\quad.

\begin{center}
\textsc{Gauss} an \textsc{Olbers}, Braunschweig, 4. Januar 1803.
\end{center}
\noindent\hspace*{\parindent}.\quad.\quad.\quad.\quad.\quad Mit meinen vorläufigen Tafeln der {\AstroSymbols ⚳}-Störungen .\quad.\quad.\quad.\quad habe ich schon oft gerechnet und ihre vorzügliche \textls[180]{Bequemlichkeit} bestätigt gefunden. Aber mit den VIII. im Nov. der M. C. abgedruckten Elementen bin ich noch nicht recht zufrieden. Ich hatte dafür ausschließlich Meridianbeobachtungen verwendet, die durch diese Elemente sehr gut und viel besser als durch die\,\raisebox{.65ex}{.} VII. dargestellt wurden. Doch seit die ohne Tafeln so \SourceLeaf{404}beschwerlichen Berechnungen der Störungen kein Hindernis mehr waren, habe ich auch die letzten Beobachtungen von \textsc{Oriani} damit verglichen und zu meiner Verwunderung den Fehler merklich größer gefunden als bei den VII. Elementen. Die Rektaszensionen stimmen alle ziemlich gut, aber die Deklinationen sind recht einheitlich etwa \(40''\) zu klein. Daher setze ich den Längenfehler für den 5. Aug. im Mittel mit \(+18''\) und den Breitenfehler mit \(-37''\) an. Der Fehler der VII. Elemente ist nur etwa halb so groß. Mir scheint, dies bestätigt eine Meinung, die ich schon früher geäußert habe. Die Unsicherheit infolge der noch zu kurzen Dauer der Beobachtungen kann sich vorerst auf die Vorhersage der künftigen Bewegung nachteiliger auswirken als die Vernachlässigung der Störungen. Ich will daher die Elemente noch weiter verfeinern und möglichst genau an alle vorhandenen Beobachtungen anpassen. Fast alle dafür nötigen Rechnungen müsste ich ohnehin in diesem Jahr 1803 ausführen, um die künftigen Beobachtungen auswerten zu können. Die Arbeit ist also nicht vergeblich, obwohl sie zur Wiederauffindung am östlichen Himmel nicht unbedingt nötig ist. Dafür werden, denke ich, die von \textsc{Triesnecker} und \textsc{Bode} berechneten Ephemeriden hinreichend genau sein. — Was eine hohe Genauigkeit der Elemente noch verhindert, ist vor allem die kurze Dauer der Beobachtungen \textsc{Piazzi}s von 1801. Für eine genaue Bestimmung ist es, wenn sich die mittlere Bewegung wie hier nicht aus zeitlich weit auseinanderliegenden Beobachtungen ableiten lässt, \textls[180]{wesentlich}, 4 Beobachtungen \(A,B,C,D\) zu haben, die alle weit genug auseinanderliegen. Auch sollte das Intervall \(B,C\) möglichst viel größer sein als \(A,B\) und\,\raisebox{.8ex}{.} \(C,D\). Dafür muss man als \(A\) und \(B\) die ersten und letzten Beobachtungen von \textsc{Piazzi} nehmen, die nur 41 Tage auseinanderliegen. Man kann sich also einen viel höheren Grad an Genauigkeit versprechen, wenn man von \textsc{Piazzi}s Beobachtungen nur noch eine für \(A\) braucht (die man natürlich durch Vergleich mit den benachbarten genauer macht, als eine einzelne Beobachtung ist), also 1803
\[\textstyle
\text{wo man }\left\{\begin{array}{c}A\\B\\C\\D\end{array}\right\}
\text{ aus }\left\{\begin{array}{ll}
1801&\\
1802&\text{(oder \textsc{Zach}s Beob. vom 7. Dez. 1801)}\\
1802&\\
1803&
\end{array}\right.
\]
\SourceLeaf{405}nehmen kann. Eine wirklich hohe Genauigkeit darf man aber erst nach der {\AstroSymbols ☍} von 1804 erhoffen, wenn man \(A,B,C,D\) aus ebenso vielen verschiedenen Jahren nehmen kann. 1806, wenn man 4 Oppositionen verwenden kann, wird man die Bahn der {\AstroSymbols ⚳}, vielleicht mit Ausnahme der mittleren Bewegung, hoffentlich fast so genau kennen wie die Bahnen der übrigen Planeten. Bei der {\AstroSymbols ⚴} wird alles ein Jahr länger dauern. — Finden wir aber einen der beiden Himmelskörper in einer älteren Beobachtung wieder, woran wir noch nicht verzweifeln wollen, so wird natürlich alles viel schneller gehen.\quad.\quad.\quad.\quad.\quad.

\begin{center}
\textsc{Gauss} an \textsc{Olbers}, Braunschweig, 1. März 1803.
\end{center}
\noindent\hspace*{\parindent}.\quad.\quad.\quad.\quad.\quad Meine Tafeln für die Störungen der {\AstroSymbols ⚳} hätte ich Ihnen längst geschickt, wenn von \textsc{Zach} mir nicht geschrieben hätte, dass alles im Februarheft der M. C. erscheinen würde. Da der Platz darin wohl nicht ausgereicht hat, erhalten Sie sie mit meinem nächsten Brief. Ich habe mir auch die übrigen für die Bewegung der {\AstroSymbols ⚳} nötigen Tafeln ohne großen Aufwand nur auf Minuten genau berechnet, eigentlich nur zum eigenen Gebrauch, um die erst kürzlich erhaltene Histoire Céleste durchzusehen. Ich habe sie bereits durchblättert und einige flüchtige Versuche angestellt. Ich verzweifle aber fast daran, die {\AstroSymbols ⚳} darin zu finden, obwohl ich auf mehrere Tage gestoßen bin, an denen der Unterschied der Deklination nur einige Grade betrug. Auch diese Tafel schicke ich Ihnen demnächst. Beide Tafeln sind jedoch nur vorläufig. Nach Abschluss der diesjährigen Beobachtungen werden wir viel genauere Elemente haben, und dann wird es sich lohnen, die Störungen genauer und vollständiger zu berechnen und in die Tafeln aufzunehmen.\quad.\quad.\quad.\quad.\quad.

\begin{center}
\textsc{Gauss} an \textsc{Olbers}, Braunschweig, 8. April 1803.
\end{center}
\noindent\hspace*{\parindent}.\quad.\quad.\quad.\quad.\quad Hiermit schicke ich Ihnen die seit Langem versprochenen kleinen {\AstroSymbols ⚳}- und {\AstroSymbols ⚴}-Tafeln. Die eine noch ausstehende für die {\AstroSymbols ⚴} erhalten Sie demnächst.\quad.\quad.\quad.\quad.\quad.

Erst in diesem Jahr werde ich zur weiteren Verbesserung der {\AstroSymbols ⚳}-Elemente die Methode anwenden können, die dann auch für die fortlaufende Verfeinerung brauchbar sein wird. Bisher war es dafür zu früh. Da diese \SourceLeaf{406}Methode meiner Meinung nach ein wichtiger Teil des Ganzen ist, will ich mich der Bearbeitung des Ganzen zuwenden, wenn ich diesen Teil erst \textls[180]{praktisch} erprobt und, wie das gewöhnlich geschieht, in der Praxis vervollkommnet habe.\quad.\quad.\quad.\quad.\quad.

\begin{center}
\textsc{Gauss} an \textsc{Olbers}, Braunschweig, 25. Januar 1805.
\end{center}
\noindent\hspace*{\parindent}.\quad.\quad.\quad.\quad.\quad Meine neuen Elemente der Ceres .\quad.\quad.\quad.\quad.\quad sind die folgenden.
\begin{center}
\begingroup\small\setlength{\tabcolsep}{5pt}\renewcommand{\arraystretch}{1.25}
\begin{tabular}{r r r r}
 & Epoche & Aphel & Exzentrizität\\
1801 & \(77^0\ 17'\ \ 0''{,}1\) & \(326^0\ 19'\ 59''\) & \(0{,}078\ 4929\)\\
1802 & \(155\ \ 27\ \ 34{,}2\) & \(22\ \ \ 0\) & \(4871\)\\
1803 & \(233\ \ 38\ \ \ 8{,}3\) & \(24\ \ \ 2\) & \(4814\)\\
1804 & \(312\ \ \ 1\ \ 33{,}5\) & \(26\ \ \ 3\) & \(4757\)\\
1805 & \(30\ \ 12\ \ \ 7{,}7\) & \(28\ \ \ 4\) & \(4700\)
\end{tabular}\endgroup
\end{center}
\begin{center}
\begingroup\footnotesize\setlength{\tabcolsep}{2pt}\renewcommand{\arraystretch}{1.25}
\begin{tabular}{l@{\quad}r}
Logarithmus der großen Halbachse .\ .\ .\ .\ . & \(0{,}442\ 0004\)\\
Tägliche tropische Bewegung .\ .\ .\ .\ .\ .\ .\ .\ .\ . & \(771''{,}0524\)\\
\(\AscendingNode\) 1801 .\ .\ .\ .\ .\ .\ .\ .\ .\ .\ .\ .\ .\ .\ .\ .\ .\ .\ .\ .\ .\ .\ .\ .\ .\ .\ .\ . & \(80^0\ 54'\ 46''\)\\
Neigung 1801 .\ .\ .\ .\ .\ .\ .\ .\ .\ .\ .\ .\ .\ .\ .\ .\ .\ .\ .\ .\ .\ .\ . & \(10^0\ 38'\ 13''\).
\end{tabular}\endgroup
\end{center}

Diese Elemente, die noch mit meinen alten Störungstafeln verbunden werden müssen, beruhen auf den 3 {\AstroSymbols ☍} und den Beobachtungen \textsc{Piazzi}s von 1801. Allerdings lassen sich die \textls[180]{Breiten} in den 4 Jahren mit der gefundenen Säkularänderung von \(\AscendingNode\) und Neigung nicht darstellen (das Erstere bemerkte ich schon im vorigen Jahr). Um die Beobachtungen mit meinen Störungstafeln in Einklang zu bringen, muss ich
\begin{center}
\begin{tabular}{l l l}
dem Knoten & \(0''{,}243\) & tägl. Verschiebung\\
der Neigung & \(0''{,}025\) & tägl. Abnahme
\end{tabular}
\end{center}
zuweisen. Beide Werte sind viel größer, als sich aus der Rechnung ergibt. Das liegt zweifellos nur an der Unvollständigkeit der periodischen Gleichungen für den Radiusvektor und die Breite und wird sich künftig bei einer ausführlicheren Berechnung der Störungen zeigen.
\par\medskip\begin{center}\rule{.22\linewidth}{.4pt}\end{center}
